\section{Limitations and Threats to Validity}
\label{sec:limitations-ra}

The construction inherits the parent paper's limitations and adds
limitations specific to the response side. The assumption ledger
below is restricted to the new obligations; the parent paper's ledger
covers the cryptographic, library, operational, and legal assumptions
common to both constructions.

\begin{table*}[t]
\centering
\caption{Assumption ledger for the response-side construction. The
parent paper's ledger applies in full; the rows below name the
additional substrate-level obligations.}
\label{tab:assumptions-ra}
\small
\begin{tabularx}{\textwidth}{@{}>{\raggedright\arraybackslash}p{.28\textwidth}>{\raggedright\arraybackslash}p{.34\textwidth}>{\raggedright\arraybackslash}X@{}}
\toprule
Assumption & Used for & Residual risk \\
\midrule
Substrate counter is monotonic and faithfully advanced by a real-time scheduler. & TTL expiry, auto-reversion, partition-window bounds. & A compromised counter falsifies the TTL claim. The substrate does not trust the endpoint's wall clock; it trusts the counter's monotonicity. \\
Inverse executors are content-hash-checked against the forward effect's recorded artifact. & Rollback receipt admissibility. & A forged artifact that matches the forward hash collapses the inverse check. \\
Ledger append happens before the side effect's external visibility. & Crash safety on the rollback path. & The deployment instance appends the execution receipt AFTER the side effect runs; a crash between effect and append leaves a live side effect with no execution record. \\
Bilateral cosignature on the destructive subclass is drawn from independent principals. & The bilateral reduction. & Inherited from the parent paper; collapses to one-of-one when a single actor operates both keys. \\
\bottomrule
\end{tabularx}
\end{table*}

The remaining limits are:

\begin{itemize}
  \item \textbf{TTL auto-expiry scheduler.} The substrate model
    assumes a scheduler advances the counter and triggers expiry
    callbacks. The deployment instance exposes the
    \codepath{/expire} HTTP endpoint but no background task invokes
    it. The headline composition theorem is unfalsifiable against the
    current binary without the scheduler; "terminates within its TTL"
    is operator-dependent rather than type-level.
  \item \textbf{Missing inverse executors on the destructive subclass.}
    Process-tree terminate, grant revoke, and evidence collection are
    destructive: each has a forward executor but no inverse. The
    bilateral reduction applies as an admission obligation; an
    empirical measurement of the destructive chain requires inverse
    executors that do not exist.
  \item \textbf{Crash safety on the ledger append.} The deployment
    instance runs the OS executor before the ledger append. A crash
    between effect and append leaves a live side effect with no
    execution record. A write-ahead ledger pattern (append a pending
    record before the effect, then mark succeeded after) would close
    this; the substrate model and the headline theorem do not require
    it, but the empirical chapter's "rollback closes the chain" claim
    is weakened by the crash window.
  \item \textbf{Bilateral cosignature not wired into the response
    engine.} The substrate's bilateral DSSE verifier exists; the
    response engine signs every receipt under a single endpoint
    keypair. The destructive subclass's bilateral admission is
    therefore a substrate-level obligation, not a deployed
    measurement. Wiring requires an operator-polity key custody
    design.
  \item \textbf{Single-actor bilateral collapse.} Inherited from the
    parent paper. Two keys held by a single actor satisfy two-key
    DSSE but not party-independence; the destructive admission
    reduces to one-of-one signatures. The construction here narrows
    the consequence (the destructive class becomes operator-key-
    dependent) without solving the underlying problem.
  \item \textbf{Sensor stub on macOS Endpoint Security.} The
    deployment instance declares the entitlement but does not
    subscribe via \codepath{es_new_client}; sensor-state attestation
    on response receipts is bounded to the Network Extension verdict
    path, the tool-preflight admission hook, and the package-manager
    runtime guard. The substrate model treats sensor state as an
    audited input; the deployment instance's input fidelity on the
    Endpoint Security path is below the substrate model's assumption.
  \item \textbf{Sensor-state fidelity on rollback and acknowledgement
    receipts.} The deployment instance attaches a placeholder sensor
    state on rollback and acknowledgement emitters rather than a
    fresh snapshot; a constitution that requires sensor-attested
    rollbacks sees synthetic data. The substrate model treats this
    as a deployment refinement obligation.
  \item \textbf{Rollback receipts under concurrent invocation.} The
    deployment instance has no per-execution rollback lock; manual
    rollback and TTL-driven rollback can fire concurrently. The
    egress ledger mutex serializes egress rollbacks; the file-rename
    rollbacks rely on filesystem-level atomicity of
    \codepath{rename(2)} plus a TOCTOU-racy existence check. The
    substrate model does not address concurrency.
  \item \textbf{Chain-of-rollback semantics.} A rollback receipt
    inverting another rollback receipt is not modeled. The
    \codepath{RollbackReceipt} structure has a single
    \codepath{rollbackOf} field; an inductive chain requires a
    structure the substrate does not introduce.
  \item \textbf{Bounded paper model.} The model encodes the TTL-
    bounded amendment chain composition theorem and four supporting
    bridges; it does not verify the whole deployment instance. The
    model is paper-local: it is checked by hand against the substrate's
    syntactic treaty model and is not in the substrate's proof manifest
    or theorem inventory. The deployment instance discharges the
    substrate-level obligations by code citation rather than by
    mechanized refinement.
  \item \textbf{rfl-gate on the headline theorem.} The headline
    composition theorem \thm{ttl_bounded_amendment_chain_preserves_baseline}
    does not discharge to \emph{rfl}: the paper-local mechanization
    proves it pointwise per amendment by a case split on the TTL window,
    composing the two refinement witnesses on the explicit finite
    admission domain in the in-window arm. The result is finite-domain:
    it holds on the domain the witnesses were checked against and says
    nothing about admission views outside it.
  \item \textbf{Operator-key custody at scale.} The bilateral
    admission on the destructive subclass requires an operator-polity
    key whose custody, rotation, revocation, and compromise recovery
    is outside the present construction. The parent paper's
    inheritance applies: per-kernel Ed25519 signing assumes a key-
    management story (HSM, KMS, TEE-bound key handles) the present
    construction does not specify.
  \item \textbf{Override path for false-rollback admission.} The
    substrate denies a rollback receipt whose canonical body fails
    the polity's predicates. The construction does not define which
    capability authorizes an admit-under-protest rollback, what
    crisis artifact the override produces, or how the override is
    audited. The same gap applies to the parent paper's override
    authority limitation.
  \item \textbf{Real-time semantics.} The substrate treats time as a
    monotonic counter; partition behavior across a TTL boundary is
    not modeled. The auto-reversion at TTL expiry during a partition
    is a deployment-side obligation rather than a substrate-level
    theorem.
\end{itemize}
